% !TeX root = main.tex

% ============================================================
% thesis-sync entry: 2026-07-06_materials-and-bonding
% Type: numeric data + design decision
% Insertable section block (no preamble). Requires: multirow (optional), no figures.
% Elastosil M4601 values transcribed from the Wacker datasheet image supplied
% by the author (Properties Cured, vulcanized 24 h at 23 C). EcoFlex 00-30 Yeoh
% coefficients are those used in the author's own FEM work (inchworm paper,
% Table 1). Cells marked "n.r." were not recorded in this session.
% ============================================================

\section{MATERIALS AND BONDING}
\label{sec:lizard-materials}

\subsection{Elastomers}

Two silicone elastomers are combined in the robot. The actuator bodies
(body, tail, and leg tubes) are cast in EcoFlex 00-30 (Smooth-On), a
translucent platinum-cure silicone. The end caps, the foot interfaces,
and the suction cups are cast in Elastosil M4601 (Wacker Chemie), a
reddish-brown two-component RTV-2 silicone that is substantially stiffer
and tougher than EcoFlex 00-30; in the photographs of the robot, every
reddish part is Elastosil M4601 and every translucent part is
EcoFlex 00-30. The division of roles follows the mechanical demand: the
actuator walls must be highly extensible at low pressure, whereas the
caps and suction cups must hold their shape, seal against the pneumatic
fittings, and survive repeated attachment cycles. The cured properties of Ecoflex 00-30 and
Elastosil M4601, transcribed from the manufacturer's datasheet, are
reported in Table~\ref{tab:lizard-materials}.

\begin{table}[tb]
  \caption{Properties of the two silicone elastomers used in the robot.
  Elastosil M4601 values are the manufacturer's cured properties
  (vulcanized 24~h at 23$^{\circ}$C); EcoFlex 00-30 Yeoh coefficients are
  those used in the FEM optimization of the parent actuator. Cells marked
  n.r.\ were not recorded and are to be completed from the respective
  datasheets.}
  \label{tab:lizard-materials}
  \centering
  \begin{tabular}{|l|c|c|}
    \hline
    Property & EcoFlex 00-30 & Elastosil M4601 \\ \hline
    Manufacturer & Smooth-On & Wacker Chemie \\ \hline
    Color & translucent & reddish brown \\ \hline
    Density & 1.07~g/cm$^3$ & 1.13~g/cm$^3$ \\ \hline
    Hardness & Shore 00-30 & Shore A 28 (ISO 868) \\ \hline
    Tensile strength & 1.38~N/mm$^2$ & 6.5~N/mm$^2$ \\ \hline
    Elongation at break & 900\% & 700\% \\ \hline
    Tear strength & 6.65~N/mm & $>$30~N/mm  \\ \hline
    Linear shrinkage & $<$0.1\% ($<$0.001~in./in.) & $<$0.1\% \\ \hline
  \end{tabular}
\end{table}

\subsection{Curing and Bonding Procedure}

All silicone parts are oven-cured at elevated temperature to accelerate
vulcanization: EcoFlex 00-30 parts are baked at 90$^{\circ}$C for
90~minutes, and Elastosil M4601 parts at 90$^{\circ}$C for 60~minutes.
Bonding between cured Elastosil and cured EcoFlex parts --- required at
every cap-to-tube and foot-to-leg interface --- is performed with KE44-T
silicone adhesive (Shin-Etsu Chemical); the bonded assembly is baked at
90$^{\circ}$C for 90~minutes to cure the adhesive.
